% GENERATED FILE. Edit canonical lessons, then run npm run generate:books.
% canonical-source-hash: fnv1a64:3f8014329070fe39
% canonical-lessons: MW-C126-koylo, MW-C126-chhano, MW-C126-boro, MW-C126-rassi, MW-C126-khuti

\chapter{Coal, Dung cake, Sack, Rope, Peg}
\label{ch:mw-coal-dung-cake-sack-rope-peg}
\hlchaptermodality{126}

\begin{chapteropening}
\textbf{By the end of this chapter:} \emph{I can say coal, a dung cake, a sack, a rope and a peg.}

Complete the last lesson of chapter 126: I can say coal, a dung cake, a sack, a rope and a peg.
\end{chapteropening}

\section[\texorpdfstring{\mw{कोयलो} (koylo)}{कोयलो (koylo)}]{\mw{कोयलो} (koylo) — coal}

\label{lesson:MW-C126-koylo}



\begin{quote}
Before the new one: say the Marwadi for matches, then the Marwadi for wood.
\end{quote}

\subsection*{You'll want to know: \mw{कोयलो}}
\textbf{\mw{कोयलो}} — \emph{koylo} — \textquotedblleft{}coal\textquotedblright{}.

\begin{cousinweb}
\textbf{\mw{कोयलो}} (\emph{koylo}) is a tadbhava, a word inherited from Sanskrit and reshaped by sound change. It comes from \emph{kokila}, \textquotedblleft{}lighted coal, charcoal\textquotedblright{}. CDIAL lists Marwari \emph{koilo} under it (CDIAL 3484). CDIAL takes the Sanskrit word from Middle Indo-Aryan \emph{kōila}, itself from Proto-Munda.
\end{cousinweb}

\begin{grammarlens}[title={What you've built}]
One more word for this chapter.
\end{grammarlens}

\subsection*{Your turn}
\begin{itemize}
\raggedright
  \item \emph{Say it:} \emph{koylo}
  \item \emph{Say it:} \emph{koylo}, once more
  \item \emph{Say it:} say \emph{koylo}
  \item \emph{Recall:} say the Marwadi for matches, then the Marwadi for wood, then say \emph{koylo} again
\end{itemize}

\begin{tcolorbox}[breakable,colback=teal!4,colframe=teal!35!black,title={Before you move on}]
What does \mw{कोयलो} mean? (\textquotedblleft{}Coal\textquotedblright{}.) Say it once more.
\end{tcolorbox}
\section[\texorpdfstring{\mw{छाणो} (chhāṇo)}{छाणो (chhāṇo)}]{\mw{छाणो} (chhāṇo) — a dung cake}

\label{lesson:MW-C126-chhano}



\begin{quote}
Before the new one: say the Marwadi for wood, then the Marwadi for coal.
\end{quote}

\subsection*{You'll want to know: \mw{छाणो}}
\textbf{\mw{छाणो}} — \emph{chhāṇo} — \textquotedblleft{}a dung cake\textquotedblright{}.

\begin{cousinweb}
\textbf{\mw{छाणो}} (\emph{chhāṇo}) is a tadbhava, a word inherited and reshaped by sound change. It comes from \emph{*chakana}, \textquotedblleft{}dung\textquotedblright{}, a form CDIAL reconstructs from its descendants. CDIAL gives no Marwari form in this entry (CDIAL 4952); the match is through Gujarati \emph{chāṇ}. CDIAL sets it beside \emph{*śakana}, \textquotedblleft{}dung\textquotedblright{}.
\end{cousinweb}

\begin{grammarlens}[title={What you've built}]
One more word for this chapter.
\end{grammarlens}

\subsection*{Your turn}
\begin{itemize}
\raggedright
  \item \emph{Say it:} \emph{chhāṇo}
  \item \emph{Say it:} \emph{chhāṇo}, once more
  \item \emph{Say it:} say \emph{chhāṇo}
  \item \emph{Recall:} say the Marwadi for wood, then the Marwadi for coal, then say \emph{chhāṇo} again
\end{itemize}

\begin{tcolorbox}[breakable,colback=teal!4,colframe=teal!35!black,title={Before you move on}]
What does \mw{छाणो} mean? (\textquotedblleft{}A dung cake\textquotedblright{}.) Say it once more.
\end{tcolorbox}
\section[\texorpdfstring{\mw{बोरो} (boro)}{बोरो (boro)}]{\mw{बोरो} (boro) — a sack}

\label{lesson:MW-C126-boro}



\begin{quote}
Before the new one: say the Marwadi for coal, then the Marwadi for a dung cake.
\end{quote}

\subsection*{You'll want to know: \mw{बोरो}}
\textbf{\mw{बोरो}} — \emph{boro} — \textquotedblleft{}a sack\textquotedblright{}.

\begin{cousinweb}
\textbf{\mw{बोरो}} (\emph{boro}) is a tadbhava, a word inherited and reshaped by sound change. It comes from \emph{*bora}, \textquotedblleft{}sacking, sack\textquotedblright{}, a form CDIAL reconstructs from its descendants. CDIAL gives no Marwari form in this entry (CDIAL 9320); the match is through Hindi \emph{borā}. CDIAL thinks it first meant cloth of goat's hair, and asks whether it goes back to Sanskrit \emph{bakura}, a goatskin bag.
\end{cousinweb}

\begin{grammarlens}[title={What you've built}]
One more word for this chapter.
\end{grammarlens}

\subsection*{Your turn}
\begin{itemize}
\raggedright
  \item \emph{Say it:} \emph{boro}
  \item \emph{Say it:} \emph{boro}, once more
  \item \emph{Say it:} say \emph{boro}
  \item \emph{Recall:} say the Marwadi for coal, then the Marwadi for a dung cake, then say \emph{boro} again
\end{itemize}

\begin{tcolorbox}[breakable,colback=teal!4,colframe=teal!35!black,title={Before you move on}]
What does \mw{बोरो} mean? (\textquotedblleft{}A sack\textquotedblright{}.) Say it once more.
\end{tcolorbox}
\section[\texorpdfstring{\mw{रस्सी} (rassī)}{रस्सी (rassī)}]{\mw{रस्सी} (rassī) — a rope}

\label{lesson:MW-C126-rassi}



\begin{quote}
Before the new one: say the Marwadi for a dung cake, then the Marwadi for a sack.
\end{quote}

\subsection*{You'll want to know: \mw{रस्सी}}
\textbf{\mw{रस्सी}} — \emph{rassī} — \textquotedblleft{}a rope\textquotedblright{}.

\begin{cousinweb}
\textbf{\mw{रस्सी}} (\emph{rassī}) is a tadbhava, a word inherited from Sanskrit and reshaped by sound change. It comes from \emph{raśmi}, \textquotedblleft{}rope, bridle, ray of light\textquotedblright{}. CDIAL gives no Marwari form in this entry (CDIAL 10648); the match is through Hindi \emph{rassī} and Gujarati \emph{rasī}. CDIAL marks the doubled \emph{ss} forms as taken from Punjabi.
\end{cousinweb}

\begin{grammarlens}[title={What you've built}]
One more word for this chapter.
\end{grammarlens}

\subsection*{Your turn}
\begin{itemize}
\raggedright
  \item \emph{Say it:} \emph{rassī}
  \item \emph{Say it:} \emph{rassī}, once more
  \item \emph{Say it:} say \emph{rassī}
  \item \emph{Recall:} say the Marwadi for a dung cake, then the Marwadi for a sack, then say \emph{rassī} again
\end{itemize}

\begin{tcolorbox}[breakable,colback=teal!4,colframe=teal!35!black,title={Before you move on}]
What does \mw{रस्सी} mean? (\textquotedblleft{}A rope\textquotedblright{}.) Say it once more.
\end{tcolorbox}
\section[\texorpdfstring{\mw{खूंटी} (khū̃ṭī)}{खूंटी (khū̃ṭī)}]{\mw{खूंटी} (khū̃ṭī) — a peg}

\label{lesson:MW-C126-khuti}



\begin{quote}
Before the new one: say the Marwadi for a sack, then the Marwadi for a rope.
\end{quote}

\subsection*{You'll want to know: \mw{खूंटी}}
\textbf{\mw{खूंटी}} — \emph{khū̃ṭī} — \textquotedblleft{}a peg\textquotedblright{}.

\begin{cousinweb}
\textbf{\mw{खूंटी}} (\emph{khū̃ṭī}) is a tadbhava, a word inherited and reshaped by sound change. It comes from \emph{*khuṭṭa}, \textquotedblleft{}peg, post\textquotedblright{}, a form CDIAL reconstructs from its descendants. CDIAL lists Marwari \emph{khuṭī} under it (CDIAL 3893).
\end{cousinweb}

\begin{grammarlens}[title={What you've built}]
One more word for this chapter.
\end{grammarlens}

\subsection*{Your turn}
\begin{itemize}
\raggedright
  \item \emph{Say it:} \emph{khū̃ṭī}
  \item \emph{Say it:} \emph{khū̃ṭī}, once more
  \item \emph{Say it:} say \emph{khū̃ṭī}
  \item \emph{Recall:} say the Marwadi for a sack, then the Marwadi for a rope, then say \emph{khū̃ṭī} again
\end{itemize}

\begin{tcolorbox}[breakable,colback=teal!4,colframe=teal!35!black,title={Before you move on}]
What does \mw{खूंटी} mean? (\textquotedblleft{}A peg\textquotedblright{}.) Say it once more.
\end{tcolorbox}
